\chapter{Central Venous Catheter Placement}

\section*{Overview}
A central venous catheter (CVC) is a catheter placed into a large vein (internal jugular, subclavian, or femoral) to administer fluids, medications, or monitor hemodynamics.

\section*{Alternative Names}
Central line insertion, CVC placement, Triple lumen catheter

\section*{Principle}
A long catheter is inserted into a large vein such as the internal jugular, subclavian, or femoral, and advanced so its tip lies in the superior vena cava near the right atrium. The high blood flow at this site permits reliable delivery of drugs and nutrition, pressure monitoring, and blood sampling.
\section*{History}
Werner Forssmann inserted a catheter into his own heart in 1929. Sven-Ivar Seldinger introduced the guidewire technique (Seldinger technique) in 1953.

\section*{Why It Is Done}
Ordered for administration of vasopressors, hypertonic solutions, total parenteral nutrition (TPN), long-term chemotherapy, or hemodialysis access.

\section*{Is This a Primary or Secondary Test or Procedure}
Primary vascular access method for hemodynamically unstable patients requiring vasoactive infusions.

\section*{How It Is Performed}
Under sterile conditions and ultrasound guidance. The skin is numbed, the target vein is punctured with a needle, a guidewire is advanced, and the catheter is threaded over the wire.

\section*{Preparation}
Verify coagulation profile (INR and platelets). Obtain sterile gown, gloves, mask, and catheter kit.

\section*{Contraindications and Precautions}
Active infection at the insertion site, or thrombosis of the target vein.

\section*{Risks}
Pneumothorax (especially subclavian), arterial puncture, air embolism, cardiac arrhythmias (if wire enters the atrium), or catheter-related bloodstream infection (CRBSI).

\section*{Aftercare and Recovery}
Obtain a chest X-ray immediately to confirm correct catheter tip position (cavoatrial junction) and rule out pneumothorax before using the line.

\section*{Outcome and Follow-up}
 catheter is patent, flushes easily, draws blood, and the tip is correctly positioned.

\section*{Expected Results}
Catheter tip located in the lower third of the superior vena cava; no pneumothorax; patent lumens.

\section*{Follow-up}
Daily site inspection for signs of infection; immediate removal when no longer clinically necessary.

\section*{References}
\begin{enumerate}
  \item MedlinePlus. Central Venous Catheter Placement. U.S. National Library of Medicine; 2025.
  \item Current Medical Diagnosis and Treatment. McGraw-Hill; 2025.
\end{enumerate}

\clearpage
